\subsection{Datos}
\label{sec:datos}

Los datos provienen de la Tabla~6 de Lakdawala (2025) y corresponden a las 80 corridas descritas en la introducción. La variable respuesta es la eficiencia hidráulica, expresada como proporción (entre 0 y 1). La Tabla~\ref{tab:variables} resume las variables y el nombre con el que aparecen en la base de datos.

\begin{table}[H]
\centering
\caption{Variables del estudio.}
\label{tab:variables}
\begin{tabular}{llll}
\toprule
Variable & Nombre en R & Tipo & Niveles \\
\midrule
Eficiencia hidráulica ($y$)            & \texttt{eficiencia} & Respuesta numérica   & 0.6025 a 0.7165 \\
Ángulo de la boquilla ($\alpha$)       & \texttt{angulo}     & Numérica (grados)    & 90, 92.5, 95 \\
Diámetro de la boquilla ($d_j$)        & \texttt{diametro}   & Numérica (mm)        & 10, 12, 14 \\
Distancia boquilla-rodete ($N_{sod}$)  & \texttt{distancia}  & Numérica (mm)        & 40, 50, 60 \\
Tipo de rodete                         & \texttt{rodete}     & Categórica           & Pelton (PTR), Turgo (TTR) \\
Número de boquillas ($n$)              & \texttt{boquillas}  & Categórica           & 1, 2 \\
\bottomrule
\end{tabular}
\end{table}

Los datos y el código del análisis están disponibles en \enlacerepositorio. El script \texttt{anal\_descrip.R} lee los datos desde el repositorio y produce todas las tablas, cifras y figuras de la Sección~\ref{sec:resultados}.
